\section*{Chapitre \ref{ch:b3:exoplanets} --- Exoplanètes}

\begin{solution}{exo:b3:exoplanets:1}
(a)~$K = (2\pi G/P)^{1/3}m_p/\bigl((M_\odot + m_p)^{2/3}\sqrt{1 - e^2}\bigr)$.
Jupiter~: $P = \qty{3.742e8}{s}$, $(2\pi G/P)^{1/3} = \num{1.039e-6}$ en unités
SI, $K = \num{1.039e-6}\times\num{1.898e27}/(\num{1.990e30})^{2/3}/0.9988
= \qty{12.5}{m.s^{-1}}$. La Terre~: $P = \qty{3.155e7}{s}$, $K =
\qty{0.089}{m.s^{-1}}$, soit environ \qty{9}{cm.s^{-1}}. (b)~\qty{3.5}{km/h} $=
\qty{0.97}{m.s^{-1}}$ et \qty{0.35}{km/h} $= \qty{0.097}{m.s^{-1}}$.
Le signal de Jupiter vaut treize fois la première précision~; celui de la
Terre reste en deçà même de la seconde, et le bruit propre de l'étoile s'y ajoute.
\end{solution}

\begin{solution}{exo:b3:exoplanets:2}
(a)~$(71492/695700)^2 = \num{0.0106}$, une baisse de \qty{1.06}{\percent}~;
$(6378/695700)^2 = \num{8.4e-5}$, soit 84 parties par million. (b)~$R_\star/a$ :
pour la Terre $695700/\num{1.496e8} = \num{4.7e-3}$, soit une chance sur 215~;
pour Jupiter $695700/\num{7.785e8} = \num{8.9e-4}$, une chance sur 1100.
\end{solution}

\begin{solution}{exo:b3:exoplanets:3}
$m_p\sin i \approx K M_\odot^{2/3}(P/2\pi G)^{1/3}$ avec $P =
\qty{3.655e5}{s}$ : $(P/2\pi G)^{1/3} = \num{9.55e4}$ en unités SI, donc $m_p
\sin i = 59\times(\num{1.9884e30})^{2/3}\times\num{9.55e4} =
\qty{8.9e26}{kg}$, soit 0.47 masse de Jupiter ou 149 masses terrestres (la
correction pour $m_p$ dans $M_\star + m_p$ est inférieure à un millième).
$a = \bigl(GM_\odot P^2/4\pi^2\bigr)^{1/3} = \qty{7.66e9}{m} =
\qty{0.051}{au}$.
\end{solution}

\begin{solution}{exo:b3:exoplanets:4}
(a)~$4GM/c^2 = 4\times\num{6.674e-11}\times\num{9.94e29}/(\num{2.998e8})^2 =
\qty{2953}{m}$ ; $(D_S - D_L)/(D_LD_S) = 1/D_S\times(D_S - D_L)/D_L =
1/(\qty{8}{kpc}) = 1/\qty{2.469e20}{m}$. $\theta_\mathrm E =
\sqrt{2953/\num{2.469e20}} = \qty{3.46e-9}{rad} = \qty{0.71}{mas}$. Au
niveau de la lentille, $R_\mathrm E = \theta_\mathrm E D_L = \num{3.46e-9}\times
\qty{1.234e20}{m} = \qty{4.27e11}{m} = \qty{2.85}{au}$. (b)~$A(1) =
3/\sqrt5 = 1.34$ ; $A(0.5) = 2.25/(0.5\times2.062) = 2.18$ ; $A(0.1) =
2.01/(0.1\times2.002) = 10.0$. (c)~$t_\mathrm E = \num{4.27e11}/\num{2e5}
= \qty{2.13e6}{s}$, soit environ 25 jours~: les événements de microlentille durent des semaines.
\end{solution}

\begin{solution}{exo:b3:exoplanets:5}
(a)~$a = (GM_\odot P^2/4\pi^2)^{1/3} = \qty{6.09e9}{m}$ (\qty{0.0407}{au}),
donc $a/R_\star = 8.75$. (b)~Avec $P/\pi = \qty{22.9}{h}$, pour $b = 0$ :
\[
T_{14} = \frac P\pi\arcsin\frac{1.12}{8.75} = \qty{2.94}{h}, \qquad
T_{23} = \frac P\pi\arcsin\frac{0.88}{8.75} = \qty{2.31}{h}.
\]
Pour $b = 0.7$ : $\cos i = 0.7/8.75$, $i = \ang{85.4}$, $(a/R_\star)\sin i =
8.726$, et
\[
T_{14} = \frac P\pi\arcsin\frac{\sqrt{1.2544 - 0.49}}{8.726} = \qty{2.30}{h},
\qquad
T_{23} = \frac P\pi\arcsin\frac{\sqrt{0.7744 - 0.49}}{8.726} = \qty{1.40}{h}.
\]
(c)~L'entrée et la sortie occupent $(T_{14} - T_{23})/T_{14} = \qty{21}{\percent}$
du transit lorsqu'il est central, \qty{39}{\percent} pour $b = 0.7$ ; un
transit rasant, $b \to 1$, n'a pas de fond plat et ressemble à un V.
\end{solution}

\begin{solution}{exo:b3:exoplanets:6}
$\bar\rho = M/\tfrac43\pi R^3$, avec $\tfrac43\pi R_\oplus^3 =
\qty{1.087e21}{m^3}$ : TRAPPIST-1~b $1.374\times\num{5.97e24}/(1.116^3\times
\num{1.087e21}) = \qty{5.43}{g.cm^{-3}}$ ; Kepler-10~b
\qty{5.60}{g.cm^{-3}}~; WASP-39~b \qty{0.167}{g.cm^{-3}}~; HD~209458~b
\qty{0.337}{g.cm^{-3}}. Les deux premières sont aussi denses que la Terre
(\qty{5.49}{g.cm^{-3}} avec les mêmes constantes) et rocheuses~; les deux
dernières sont des planètes géantes gazeuses, plus légères que l'eau, WASP-39~b
étant huit fois moins dense que Jupiter.
\end{solution}

\begin{solution}{exo:b3:exoplanets:7}
Les extrema valent $\gamma + K(1 + e\cos\omega)$ et $\gamma - K(1 -
e\cos\omega)$. Leur différence vaut $2K = \qty{60.0}{m.s^{-1}}$, donc $K =
\qty{30.0}{m.s^{-1}}$ ; leur moyenne vaut $\gamma + Ke\cos\omega =
\qty{20.4}{m.s^{-1}}$, donc $Ke\cos\omega = \qty{8.0}{m.s^{-1}}$ et
$e\cos\omega = 0.27$. Sur une orbite excentrique, l'étoile reste longtemps près
de l'apoastre et traverse rapidement le périastre~; la courbe est donc
dissymétrique, et le milieu de ses extrema n'est pas sa moyenne temporelle $\gamma$.
\end{solution}

\begin{solution}{exo:b3:exoplanets:8}
(a)~$g = GM/R^2 = \num{6.674e-11}\times\num{5.332e26}/(\num{9.144e7})^2 =
\qty{4.26}{m.s^{-2}}$ ; $H = kT/\bar m g = \num{1.381e-23}\times1100/(2.3
\times\num{1.661e-27}\times4.26) = \qty{934}{km}$. Avec $R_\star =
\qty{6.533e8}{m}$, la profondeur vaut $(\num{9.144e7}/\num{6.533e8})^2 =
\num{0.0196}$. (b)~$2HR_p/R_\star^2 = 2\times\num{9.34e5}\times
\num{9.144e7}/(\num{6.533e8})^2 = \num{4.0e-4}$, soit 400 parties par million~;
deux hauteurs d'échelle, 800~: une géante gazeuse de faible gravité est la cible la plus favorable
pour la \omterm{def:b3:exoplanets:transmission}{spectroscopie de transmission}.
\end{solution}

\begin{solution}{exo:b3:exoplanets:9}
$p = 1.116\times6378/(0.1192\times695700) = 0.0858$, $p^2 = \num{7.37e-3}$.
$B_\lambda \propto 1/(e^{hc/\lambda kT} - 1)$ avec $hc/\lambda k =
\qty{959}{K}$ à \qty{15}{\micro m}~: l'étoile $1/(e^{0.374} - 1) = 2.20$ ;
une face diurne à \qty{400}{K}, $1/(e^{2.40} - 1) = 0.0998$, rapport \num{0.0453},
profondeur \num{3.3e-4} (334 parties par million)~; à \qty{500}{K}, $1/(e^{1.92} -
1) = 0.172$, rapport \num{0.078}, profondeur \num{5.7e-4}. Une centaine de kelvins
modifie la profondeur de \qty{70}{\percent}~: la profondeur mesurée fixe donc
étroitement la température de la face diurne, et une face diurne chaude signifie
que peu de chaleur est transportée vers la face nocturne.
\end{solution}

\begin{solution}{exo:b3:exoplanets:10}
(a)~$S \propto L/d^2$, donc $d = \sqrt{0.005}\,\unit{au} = \qty{0.0707}{au}$ ;
$P = 2\pi\sqrt{a^3/GM} = \qty{1.33e6}{s} = \qty{15.4}{d}$.
(b)~$K = (2\pi G/P)^{1/3}M_\oplus/(0.2M_\odot)^{2/3} = \qty{0.75}{m.s^{-1}}$,
soit 8.4 fois l'effet de la Terre sur le Soleil~: l'étoile légère se déplace davantage,
et la courte période y contribue aussi. (c)~$(6378/139140)^2 = \num{2.1e-3}$, soit 25 fois
le transit de la Terre devant le Soleil. Les deux signaux sont plus forts, les transits
reviennent toutes les deux semaines, et la probabilité de transit $R_\star/a =
\num{0.013}$ est trois fois celle de la Terre~: les petites étoiles offrent aux planètes
rocheuses de leur zone habitable une chance d'être détectées et étudiées.
\end{solution}

\begin{solution}{exo:b3:exoplanets:11}
Avec $r\,dr = -\mu\,d\mu$, l'intensité moyenne vaut $\frac1\pi\int_0^1
I\,2\pi r\,dr = 2\int_0^1 I(\mu)\mu\,d\mu = 2\bigl[\tfrac12 - u_1(\tfrac12 -
\tfrac13) - u_2\int_0^1(1 - \mu)^2\mu\,d\mu\bigr]$. La dernière intégrale vaut
$\tfrac12 - \tfrac23 + \tfrac14 = \tfrac1{12}$, donc la moyenne vaut $1 - u_1/3 -
u_2/6$, soit ici $0.825$. Une petite planète masque $\pi p^2 I(\mu)$ sur $\pi\times
0.825$~: profondeur $p^2 I(\mu)/0.825$. En $r = 0$, $I = 1$~: $1.212\,p^2$~; en
$r = 0.5$ ($\mu = 0.866$), $I = 0.942$~: $1.142\,p^2$~; en $r = 0.9$ ($\mu =
0.436$), $I = 0.695$~: $0.842\,p^2$. Un modèle de disque uniforme lirait
$1.212\,p^2$ comme $p^2$ et surestimerait le rayon de $\sqrt{1.212} = 1.10$,
\qty{10}{\percent}.
\end{solution}

\begin{solution}{exo:b3:exoplanets:12}
(a)~$R_\odot/a$~: \num{9.3e-3} à \qty{0.5}{au}, \num{6.6e-3} à
\qty{0.7}{au}, \num{4.7e-3} à \qty{1}{au}. (b)~$42000\times0.11\times
\num{6.6e-3} = 31$ planètes transitent. (c)~Seul environ un tiers d'entre elles
ont été trouvées~: les autres donnaient des transits trop peu profonds (84 parties
par million pour une Terre) ou trop rares (une ou deux par an) pour émerger du
bruit. Cette fraction est l'efficacité de détection~; sans corriger ni celle-ci
ni la probabilité de transit, les comptages sous-estimeraient la fréquence de ces
planètes d'un facteur d'environ 460.
\end{solution}

\begin{solution}{pb:b3:exoplanets:1}
\textbf{1.} $M_\star = 0.80\times\num{1.9884e30} = \qty{1.591e30}{kg}$, $P =
\qty{1.0368e6}{s}$: $a = (GM_\star P^2/4\pi^2)^{1/3} = \qty{1.4245e10}{m}
= \qty{0.0952}{au}$.

\textbf{2.} $v = 2\pi a/P = \qty{86.3}{km.s^{-1}}$.

\textbf{3.} $m_p\sin i = KM_\star^{2/3}(P/2\pi G)^{1/3} = 2.59\times\num{1.363e20}
\times\num{1.352e5} = \qty{4.77e25}{kg}$, \num{7.99} masses terrestres.

\textbf{4.} $a_\star = a\,m_p/M_\star = \num{1.4245e10}\times\num{4.77e25}/
\num{1.591e30} = \qty{427}{km}$; $2\pi a_\star/P = \qty{2.59}{m.s^{-1}}$,
la valeur $K$ observée.

\textbf{5.} $2.59$ fois une précision de \qty{1}{m.s^{-1}}, 26 fois une précision de
\qty{0.1}{m.s^{-1}}~: la planète est à la portée des meilleurs instruments,
à peine de la génération précédente.

\textbf{6.} $m_p\sin i \propto K\sqrt{1 - e^2}$~: $\times0.980$, 7.83 masses
terrestres, \qty{2}{\percent} de moins.

\textbf{7.} $p = \sqrt{\num{8.38e-4}} = 0.02895$.

\textbf{8.} $R_\star = 0.76\times\qty{695700}{km} = \qty{528700}{km}$; $R_p =
\qty{15310}{km} = 2.40$ rayons terrestres.

\textbf{9.} $a/R_\star = \num{1.4245e10}/\num{5.287e8} = 26.94$.

\textbf{10.} $T_{14} = (\qty{288}{h}/\pi)\arcsin(1.02895/26.94) = \qty{3.50}{h}$.

\textbf{11.} $\sqrt{(1 + p)^2 - b^2} = (a/R_\star)\sin(\pi T_{14}/P) = 26.94\times
\sin(\pi\times3.23/288) = 0.949$ (avec $\sin i \approx 1$)~; $b^2 =
1.0587 - 0.9008 = 0.158$, $b = 0.40$.

\textbf{12.} $\cos i = b\,R_\star/a = 0.40/26.94 = 0.0148$, $i = \ang{89.15}$.

\textbf{13.} $R_\star/a = 1/26.94 = 0.037$~: une chance sur 27.

\textbf{14.} $\sin i = 0.99989$~: la masse vraie vaut 7.99 masses terrestres, la \omterm{def:b3:exoplanets:rv}{masse
minimale} à \num{1e-4} près.

\textbf{15.} $\bar\rho = \num{4.77e25}/\bigl(\tfrac43\pi(\num{1.531e7})^3\bigr) =
\num{4.77e25}/\num{1.502e22} = \qty{3180}{kg.m^{-3}}$,
\qty{3.18}{g.cm^{-3}}.

\textbf{16.} La planète est $2.40/1.74 = 1.38$ fois plus grande qu'une planète rocheuse
de sa masse, 2.6 fois son volume. La roche seule ne suffit pas~: il lui faut
une enveloppe de gaz léger, quelques pour cent de sa masse en hydrogène et en hélium,
ou une forte fraction d'eau.

\textbf{17.} À 2.40 rayons terrestres, elle se situe au-dessus de la vallée, dans le
pic de 2 à 3 rayons terrestres~: une \omterm{def:b3:exoplanets:classes}{sous-Neptune}.

\textbf{18.} $g = Gm_p/R_p^2 = \num{6.674e-11}\times\num{4.77e25}/(\num{1.531e7})^2 =
\qty{13.6}{m.s^{-2}}$.

\textbf{19.} $L = 4\pi R_\star^2\sigma T^4 = 4\pi(\num{5.287e8})^2\times\num{5.670e-8}
\times5000^4 = \qty{1.245e26}{W}$, soit 0.325 fois la luminosité solaire.

\textbf{20.} $S/S_\oplus = 0.325/0.0952^2 = 35.9$~: bien en deçà du bord interne de
la zone habitable.

\textbf{21.} $L/4\pi a^2 = \qty{4.88e4}{W.m^{-2}}$; $T_\mathrm{eq} =
\bigl(\num{4.88e4}\times0.7/(4\sigma)\bigr)^{1/4} = \qty{623}{K}$.

\textbf{22.} $H = kT/\bar m g = \num{1.381e-23}\times623/(2.3\times\num{1.661e-27}
\times13.6) = \qty{166}{km}$.

\textbf{23.} $2HR_p/R_\star^2 = 2\times\num{1.66e5}\times\num{1.531e7}/
(\num{5.287e8})^2 = \num{1.82e-5}$, 18 parties par million~; une bande de deux
hauteurs d'échelle, 36 parties par million.

\textbf{24.} $H$ est $18/2.3 = 7.8$ fois plus petit, \qty{21}{km}~: 2.3 parties par
million par hauteur d'échelle, 4.6 pour la bande. À dix parties par million,
l'atmosphère riche en hydrogène montre sa bande d'eau à plus de trois fois le bruit~;
une atmosphère de vapeur apparaît plate, comme une atmosphère nuageuse ou une planète
sans air.

\textbf{25.} \textbf{\qty[text-series-to-math=true]{3.2}{g/cm^3}}.
\end{solution}
